\section{Introduction}
\label{sec:intro}

When a large language model generates a factual answer, uncertainty manifests
as \emph{inter-generational variance}: ask the same question ten times and you
may observe different answers. Semantic entropy~\cite{kuhn2023semantic} and
related methods detect this variance by clustering multiple responses and
measuring their semantic spread. The underlying assumption --- that uncertainty
produces observable inconsistency across outputs --- holds for factual
generation.

Code generation breaks this assumption. When GPT-4o is asked to integrate
the Stripe payment API, it chooses between the deprecated \texttt{Charge.create}
method and the current \texttt{PaymentIntent} workflow. This is a consequential
decision. But the model has a strong training-data prior for one choice: it
generates the same API call across all samples, consistently and confidently.
The uncertainty that exists at generation time --- visible in the token-level
log-probability distribution as a near-50/50 split between the two API paths
--- is committed and gone from the output. Repeated sampling detects nothing.

We call this the \emph{intra- vs.\ inter-generation distinction}.
Code generation uncertainty is intra-generational: it lives inside a single
forward pass, token by token, and must be measured there.
This observation is not merely conceptual --- we demonstrate it empirically
in Section~\ref{sec:comparison}: UQLM~\cite{uqlm2024} (a multi-sample
semantic-entropy tool), the closest published tool for LLM uncertainty
quantification, achieves 5\% detection on a benchmark of API-version-split
code generation prompts where our system achieves 85\%.

The practical consequence is that AI code generation tools --- Copilot, Cursor,
Claude, ChatGPT --- produce
output that looks uniformly confident. No signal distinguishes a generation
where the model was certain from one where it was choosing between a deprecated
API that will crash in production and its modern replacement.
The consequences are not evenly distributed. A sorting function with an unusual
variable name is harmless. A payment processing function that calls a removed
API method is a production incident. A FastAPI endpoint that uses
\texttt{requests.get()} inside an \texttt{async def} will pass all unit tests
and silently block the event loop under real load. The developer who reviews
all generated code equally is reviewing the wrong things.
The intra-generational nature of this uncertainty means the only viable
measurement point is inside the generation itself --- at the log-probability
distribution over token alternatives at each position.

We introduce $\eps$ (epsilon), a token-level uncertainty primitive computed
from the log-probability distribution returned by the model during a single
generation. For each token position, $\eps_t$ is the normalized Shannon entropy
over the top-5 alternative tokens. The file- or function-level score is the
maximum $\eps_t$ across consequential code tokens.
Three design choices distinguish $\eps$ from naive entropy measurement:

\begin{enumerate}[noitemsep]
  \item \textbf{AST-based cosmetic filtering.} The model's uncertainty about
    whether to name a parameter \texttt{n} or \texttt{num} is real entropy but
    zero risk. We parse the generated code with Python's \texttt{ast} module
    and exclude tokens at declaration positions (function names, parameter
    names, local variable names) from aggregation.

  \item \textbf{Max aggregation, not compound.} The compound formula
    $1 - \prod(1 - \eps_i)$ converges to 1.0 for any non-trivial response
    length, making it useless as a threshold signal. Max asks the right
    question: \emph{how uncertain was the model at its single worst consequential
    decision?}

  \item \textbf{Zero added latency.} Log-probabilities are returned by the
    API alongside the generation. The $\eps$ computation runs on the already-
    received data; no additional API call is required.
\end{enumerate}

This paper makes four contributions:

\begin{enumerate}[noitemsep]
  \item \textbf{Intra- vs.\ inter-generation distinction} (\S\ref{sec:comparison}).
    A conceptual framing that explains why multi-sample uncertainty quantification
    methods are structurally ineffective for code generation API decisions, with
    empirical confirmation: 5\% vs.\ 85\% detection rate on the same benchmark
    (UQLM vs.\ $\eps$, GPT-4o).

  \item \textbf{The $\eps$ primitive} (\S\ref{sec:epsilon}).
    A single-generation, token-level, AST-filtered uncertainty metric for code,
    with a three-tier threshold protocol and adaptive ensemble threshold.

  \item \textbf{Failure taxonomy} (\S\ref{sec:evaluation}).
    Four classes of LLM code generation failure --- deprecated API, confident
    wrongness, syntax split, and compatibility mismatch --- each demonstrated
    with a calibrated scenario and live API results.

  \item \textbf{Uncertainty taxonomy} (\S\ref{sec:epsilon}).
    Three types of token-level uncertainty in generated code: cosmetic,
    consequential, and implementation-approach. An empirical benchmark
    characterizes each type and demonstrates AST filtering's effect on
    false positive rate.
\end{enumerate}

The implementation, benchmark, and all scenario scripts are available at
\url{https://github.com/michaeljaquith/epsilon-code}.
